\documentclass[11pt]{article}
\usepackage[margin=1in]{geometry}
\usepackage{amsmath,amssymb,amsthm,booktabs,hyperref}
\hypersetup{colorlinks=true,urlcolor=blue}
\newtheorem{theorem}{Theorem}
\newtheorem{corollary}[theorem]{Corollary}
\newcommand{\Z}{\mathbb Z}
\title{Integer diagonal forms for box sums and Hardin's king-neighbor table}
\author{}
\date{}
\begin{document}
\maketitle
\begin{abstract}
The base-two logarithm of the number of binary king-neighbor equality maps on an $n\times k$ board equals $(n-\mathbf1_{n\equiv2\bmod3})(k-\mathbf1_{k\equiv2\bmod3})$. This proves the empirical formula in OEIS A219525 as a consequence of established Lights Out results. We give a uniform integer diagonalization for additive box sums with arbitrary left and right reaches on each axis and independently chosen open or periodic boundaries. It yields the number of outputs and the multiplicity of every attainable output over any finite abelian group.
\end{abstract}

\section{One axis}
Fix a length $n\ge1$ and nonnegative left and right reaches $a,b$. Set $L=a+b+1$. The matrix $B$ sends $(x_1,\ldots,x_n)$ to the vector of consecutive sums
\[
 (Bx)_i=\sum_{j=i-a}^{i+b}x_j\qquad (1\le i\le n).
\]
On an open axis, set $x_j=0$ outside $[1,n]$. On a periodic axis, take indices modulo $n$ and retain multiplicities if $L>n$.

\begin{theorem}
For an open axis, form the following subsets of $\Z/L\Z$:
\[
 U=\{1-a,\ldots,0\}\bmod L,\qquad V=\{n+1,\ldots,n+b\}\bmod L .
\]
Intervals of length zero give the empty set. With $d=|U\cap V|$, the Smith normal form of $B$ has $n-d$ ones and $d$ zeros. For $a=b=r$, this nullity simplifies to
\[
 d=\max\{0,\min(s-1,2r+1-s)\}, \qquad s=n\bmod(2r+1),\quad 0\le s<2r+1.
\]
For a periodic axis, set $g=\gcd(n,L)$. Then its Smith diagonal is
\[
 \underbrace{1,\ldots,1}_{n-g},\quad L/g,\quad
 \underbrace{0,\ldots,0}_{g-1}.
\]
The periodic case is the consecutive-ones circulant formula of Williams \cite[Corollary 3.3(i)]{williams}. We include its argument to fix the boundary and multiplicity conventions.
\end{theorem}
\begin{proof}
Present the abelian group $\Z^n$ modulo the subgroup generated by the rows of $B$. This group is the cokernel of $B^{\mathsf T}$ and determines the Smith diagonal since $B$ and $B^{\mathsf T}$ have identical Smith forms. Denote the standard generators by $e_j$. In the open case put $e_j=0$ at the $a$ left and $b$ right exterior positions. The row relations $R_i=\sum_{j=i-a}^{i+b}e_j=0$ can be replaced by $R_1=0$ and all their consecutive differences,
\[
                e_{j+L}=e_j\qquad (1-a\le j\le n-1-a).
\]
On the extended interval $[1-a,n+b]$, these are precisely the edges linking successive representatives of the same residue modulo $L$. Every residue occurs because this interval has length $n+L-1\ge L$. Thus one free generator per residue survives, except that every residue in $U\cup V$ has been set to zero. The first row runs across one complete residue system, so its relation is the sum of all surviving generators, each once. There are $L-|U\cup V|=1+d$ such generators. Removing one using this relation leaves a free abelian group of rank $d$. This proves the open Smith form.

For $a=b=r$, the two forbidden arcs have length $r$ in a cycle of length $2r+1$. Their starting points differ by $n+r$ modulo $2r+1$. If $c$ is the shorter circular distance between these points, their overlap is $r-c$. Substitution of the residue $s$ gives the formula, including $r=0$.

For a periodic axis, taking differences around the cycle identifies each $e_i$ with $e_{i+L}$. The wraparound difference follows from the others, so retaining $R_1$ and all these differences preserves exactly the original relations. The surviving generators are indexed by the $g$ residues modulo $g$. In $R_1$ each appears $L/g$ times. The sum of the $g$ generators is a primitive basis vector. Thus this group is $\Z^{g-1}\oplus\Z/(L/g)\Z$, giving the claimed diagonal.
\end{proof}

\section{Several axes and finite coefficients}
Let the axis lengths be $n_1,\ldots,n_m$. Independently choose both reaches and a boundary convention on each axis. The output at a grid point is the sum over all displacement vectors in the resulting box. Let $D_i$ be the list of the $n_i$ one-axis Smith entries in Theorem~1.

\begin{corollary}
The integer matrix of box sums is unimodularly equivalent to the diagonal list $\mathcal D$ consisting of every product $d_1\cdots d_m$ with one $d_i$ chosen from each list $D_i$, retaining multiplicities. For a finite abelian group $A$, put $A[t]=\{x\in A:tx=0\}$, including $A[0]=A$. The number of possible output arrays and the number of input arrays over each attainable output are respectively
\[
       \prod_{t\in\mathcal D}\frac{|A|}{|A[t]|},
       \qquad \prod_{t\in\mathcal D}|A[t]|.
\]
For $A=\Z/q\Z$, $|A[t]|=\gcd(q,t)$.
\end{corollary}
\begin{proof}
With lexicographic ordering the box matrix is the Kronecker product of the one-axis matrices, since their entries multiply to count each displacement's multiplicity. Tensoring the left and right unimodular changes of basis preserves invertibility over the integers, with the tensor product of the inverse changes giving the inverse. The resulting diagonal products act coordinate by coordinate. Multiplication by $t$ on $A$ has kernel $A[t]$ and an image of order $|A|/|A[t]|$. Every fiber over an attainable value is a coset of the kernel.
\end{proof}

For example, take a $5\times4$ array with radius one, open in the first coordinate and periodic in the second. Its open factor has four ones and one zero. Its periodic factor has three ones and a three. The box diagonal is therefore twelve ones, four threes and four zeros. Over $\Z/6\Z$ the image has $6^{12}2^4$ elements and each attainable output has $3^4 6^4$ preimages.

\section{The empirical OEIS formula}
In Hardin's A219525 \cite{hardin}, the input is binary. At each cell the output bit indicates whether its value equals the sum modulo two of all its king-move neighbors. This indicator equals one plus the sum of the full three-by-three window, clipped to the board. Thus the output is a fixed translate of the additive image over $\mathbb F_2$. In particular the base-two logarithm of its size is the rank of the box matrix.

\begin{corollary}
Writing $f(n)=n-\mathbf1_{n\equiv2\bmod3}$, one has
\[
                   T(n,k)=f(n)f(k)\qquad (n,k\ge1).
\]
\end{corollary}
\begin{proof}
At radius one the open-axis nullity from Theorem~1 is one at lengths congruent to two modulo three and zero otherwise. The product diagonal has exactly $f(n)f(k)$ nonzero entries. To compare with the printed formula, let $s=n\bmod3$. The identities
\[
 3\lfloor(n-1)/3\rfloor+s^2-3s+3=f(n),\qquad
 3n-3(s^2-s)/2=3f(n)
\]
follow for $s=0,1,2$. Inserting them in Hardin's expression and applying the first identity to $k$ gives precisely $f(n)f(k)$.
\end{proof}

\paragraph{Prior work.}
The open-path nullity at radius one appears in the Lights Out literature \cite{ballard}. The strong-product adjacency matrix identity gives the tensor product for its closed-neighborhood map \cite{munter}. Together with multiplicativity of field rank, these imply the binary OEIS result directly. The cyclic one-axis Smith form also precedes this note \cite{williams}. The additional formulation here combines the explicit overlap rule for arbitrary asymmetric open windows with periodic factors and gives image and fiber sizes for every finite abelian coefficient group. This work uses a conventional written proof and a separate finite verification script.

\begin{thebibliography}{9}
\bibitem{hardin} R. H. Hardin, A219525, OEIS Foundation Inc., November 21, 2012. \url{https://oeis.org/A219525}
\bibitem{ballard} L. E. Ballard, E. L. Budge and D. R. Stephenson, Lights Out for graphs related to one another by constructions, \emph{Involve} 12 (2019), 181--201. \url{https://doi.org/10.2140/involve.2019.12.181}
\bibitem{munter} R. Munter and T. Peters, LIGHTS OUT! on graph products over the ring of integers modulo $k$, \emph{Electronic Journal of Linear Algebra} 37 (2021), 416--424. \url{https://doi.org/10.13001/ela.2021.5601}
\bibitem{williams} G. Williams, Smith forms for adjacency matrices of circulant graphs, \emph{Linear Algebra and its Applications} 443 (2014), 21--33. \url{https://doi.org/10.1016/j.laa.2013.11.006}
\end{thebibliography}
\end{document}
